% 本文件由 LaTeX 排版 Agent 根据有效 translation.json 自动生成。
% author=Gregory of Nyssa; work_id=2908; title=大教理

\chapter{大教理}

% structure: p0001 source=h1 target=omitted reason=page-title-already-rendered-by-parent

\Emph{天主圣三。}\par

序言及第一章。——相信天主，基于世界秩序所展示的艺术与智慧；相信天主的\Emph{唯一性}，则基于祂在能力、美善、智慧等方面所必然具有的\Emph{完美}。然而，基督徒在反对多神论时需谨慎，免得在与希腊教争辩时，不自觉地陷入犹太教。因为天主拥有圣言：否则祂就没有理性。而这圣言不能仅仅是天主的一个属性。我们被引向一个更崇高的圣言概念，基于此：天主越比我们伟大，祂的一切属性也必然比我们自身的属性更高。我们的\OriginalTerm{言语}{logos}是有限而短暂的；但神圣圣言的实体必须是不朽的；同时是活的，因为理性的不能像石头那样无生命。它也必须拥有独立的生活，而非分有的生活，否则就失去了单纯性；作为活的，它还必须具有意志的能力。圣言的意志必须与祂的能力相称：因为选择与无能的混合又会破坏单纯性。祂的意志，作为神圣的，也必须是美善的。由这种行善的能力与意志，就产生了美善的实现；因而产生了智慧与艺术地安排的世界。然而，更进一步，圣言的概念在某种意义上是相对性的，因此，随着圣言，必须承认那言说圣言者，即圣言之父，存在。这样，信德的奥秘就同样避免了犹太一神论的荒谬和异教多神论的荒谬。一方面，我们说圣言具有\Emph{生命与行动}；另一方面，我们肯定在源自圣父的\OriginalTerm{圣言}{\GreekText{Λόγος}}中，我们发现\Emph{所有}圣父本性的属性。\par

第二章。——通过人类呼吸的类比（呼吸无非是吸入和呼出的气，即一种外在于我们的事物），证明了圣神与天主本质的共融，以及祂存在的独立性。\par

第三章。——于是，从犹太教义中保留了神圣本性的统一性；从希腊教中保留了位格的区别。\par

第四章。——从圣经中驳斥犹太人。\par

\Emph{降生成人的合理性。}\par

第五、六章。——天主藉其理性和智慧创造了世界；因为祂不可能在无理性中进行那项工作；但祂的理性和智慧，如上所示，不能被视为一种说出的言语，或仅仅拥有知识，而是一种有位格的、愿意的潜能。如果整个世界是由这第二位的神圣位格创造的，那么人当然也是由祂创造的；然而并非出于任何必然，而是\Emph{出于满溢的爱}，为了存在一个能够参与天主完美性的受造物。如果人将要接受这些完美性，那么他的本性必须包含一个与天主相似的元素；特别地，他必须是不朽的。于是，人按天主的肖像被造。因此他不可能没有自由、独立、自我决定的恩赐；他参与天主恩赐因此取决于他的德行。由于这种\Emph{自由}，他能够选择邪恶，邪恶不可能起源于天主的意愿，而只在于我们自己内部，以偏离美善的形式出现，因而成为美善的缺失。恶习与德行相对，只是作为更好的东西的缺失。既然所有受造物都处于变化之下，那么可能最初有一个受造的精神将他的目光从美善移开，并变得嫉妒，从这嫉妒产生了趋向邪恶的倾向，这倾向按照自然的顺序，为所有其他的邪恶铺平了道路。他诱惑了最初的人们，\Emph{通过扰乱天主在其感性与理性本性之间制定的和谐秩序}，使他们背离美善的愚妄；并诡诈地用邪恶玷污了他们的意志。\par

第七、八章。——天主并没有因为预知从人的创造中会产生邪恶，就不创造人；因为通过忏悔和身体痛苦的方式将罪人带回到原始的恩宠，比根本不创造人更好。扶起跌倒者是适合生命赐予者的工作，祂是天主的智慧和能力；为此目的，祂成了人。\par

第九章。——降生成人并非不配祂；因为\Emph{只有邪恶带来卑贱}。\par

第十章。——认为有限不能包含无限，因而人性不能将神性接纳到里面的异议，是基于错误的假设，即圣言的降生成人意味着天主的无限被包含在肉身的界限内，如同在容器中。——火焰与灯芯的比喻。\par

第十一、十二、十三章。——此外，神性与人性结合的方式超越了我们的理解能力；虽然我们不被允许怀疑那在耶稣身上结合的事实，\Emph{基于祂所行的奇迹}。这些奇迹的超自然性质证明了它们的神圣起源。\par

第十四、十五、十六、十七章。——降生成人的计划被进一步阐明，以显示这种人的救恩之路优于天主意志的单一命令。基督承担了人的\Emph{软弱}；但那是身体的，而非道德的软弱。换言之，神圣的美善没有变成它的对立面，对立面只是恶习。在祂内，灵魂和身体结合，然后分离，按照本性的进程；但在祂如此洁净了人类生命之后，祂\Emph{以更广泛的方式}，为所有人，并永远，在复活中重新结合了它们。\par

第十八章。——魔鬼崇拜的停止、基督徒的殉道以及耶路撒冷的荒凉，被一些人接受为降生成人的证明——\par

第十九、二十章。——但希腊人和犹太人却不如此。现在让我们回到它的\Emph{合理性}。无论我们看天主的良善、德能、智慧或公义，它都显示出所有这些公认属性的结合；若缺少其中一项，便不再是属神的。因此它符合天主的完美。\par

第二十一、二十二、二十三章。——那么其中的\Emph{公义}是什么？我们必须记得，人受造时必然是变化不定的（或变好或变坏）。道德的美应是他的自由意志所趋向的方向；但他却被那美的幻象所欺骗，走向毁灭。当我们这样\Emph{自由地}把自己卖给那欺骗者之后，那出于良善而想要恢复我们自由的那一位，不能——因为祂也是公义的——为此目的采取强暴的措施。所以必须付出赎价，其价值超过所要赎回之物；因此天主子必须将自己交于死亡之权下。于是天主的\Emph{公义}促使祂选择一种交换的方法，正如祂的\Emph{智慧}在执行中显明出来。\par

第二十四、二十五章。——但\Emph{德能}又如何？那在神性降卑为人时，比宇宙间一切自然奇观更显明地彰显出来。就如火焰向下倾注一般。然后，在如此诞生之后，基督战胜了死亡。\par

第二十六章。——确实，以人性隐藏神性，乃是向那恶者施行了一种欺骗；但对那恶者而言，他既是欺骗者，自己被欺骗乃是公正的报应：那大仇敌最终必亲自发现所行之事是公义而有益的，当他也将体验到降生成人的益处时。他，一如人性，将被净化。\par

第二十七、二十八章。——病人要得医治，必须被\Emph{触及}；而人性必须被基督触及。人性并不在天上；所以只有通过降生成人才能得医治。——此外，祂的神性取人性并不比取一个\Quote{天上的}身体更不相称；一切受造物在神性之下都是同等的。甚至\Quote{丰厚的尊荣}也当归于人生育的工具。\par

第二十九、三十、三十一章。——至于降生成人的延迟，必须等到人性败坏达到最低点，救恩工程才能介入。然而，因信德而来的恩宠并未临到所有人，这应归咎于人的自由；若天主以强力手段打破我们的反抗，人类行为的可称赞性就会被摧毁。\par

第三十二章。——甚至十字架上的死亡也是崇高的：因为在那爱的计划中，它是终极而必要的一点——死亡之后，整个人类之\Quote{团}将有福的复活；十字架本身也有奥秘的意义。\par

\Emph{圣事}。\par

第三十三、三十四、三十五、三十六章。——圣洗圣事的救恩性质取决于三件事：祈祷、水和信德。1. 它显示了祈祷如何确保天主的临在。天主是真理的天主，祂应许（正如奇迹证明祂已经来过）若以\Emph{特别的方式}呼求，祂必降临。2. 它显示了神性如何从水中赋予生命。在人的生育中，即使没有祈祷，祂也从一个微小的开端赋予生命。在更高的生育中，祂将物质转化为不是灵魂，而是精神。3. 人的自由，正如在信德和悔改中所显明的那样，对重生也是必要的。在水中浸没三次是我们最初的克己；从水中出来，预示着纯洁者将有\Emph{有福的}复活的容易：整个进程是对基督的模仿。\par

第三十七章。——圣体圣事结合身体于天主，正如圣洗圣事结合灵魂于天主。我们的身体接受了毒药，需要解药；而解药只能通过吃喝进入。一个身体，神性的容器，就是这解药，如此被接纳。但它如何完整地进入每一位信友？这需要说明。水将其自身给予皮囊。同样，营养（饼和酒）变成血肉，给人身体增加分量：营养就是身体。正如在其他人的情形中，我们救主的营养（饼和酒）曾是祂的身体；但这些营养和身体，在祂内，通过内住的圣言，被改变为天主的身体。所以现在，饼和酒，被圣言（神圣祝福）祝圣后，\Emph{同时}被改变成那圣言的身体；而这肉被分施给所有信友。\par

第三十八、三十九章。——为得重生，必须相信圣子和圣神并非受造的神体，而是与天主圣父同性同体的；因为若有人在圣洗呼求中使自己的救恩依赖于任何受造之物，他就是信赖一个不完美的本性，而这本性本身需要救主。\par

第四十章。——唯有那以去除自身一切恶习来证明自己重生的人，才真正成为天主的子女。\par

\SourceRecord{Translated by William Moore and Henry Austin Wilson.；Nicene and Post-Nicene Fathers, Second Series；Vol. 5.；Edited by Philip Schaff and Henry Wace.；Buffalo, NY: Christian Literature Publishing Co.,；1893.；<http://www.newadvent.org/fathers/2908.htm>.}
